\documentclass[11pt,a4paper]{moderncv}
\moderncvtheme[blue]{classic}
\usepackage[utf8]{inputenc}
\usepackage[inline]{enumitem}
\usepackage{xpatch}
\usepackage{lmodern}
\newcommand{\ts}{\textsuperscript}
\newcommand{\st}{\ts{st} }
\newcommand{\nd}{\ts{nd} }
\newcommand{\rd}{\ts{rd} }
\usepackage[top=3cm, bottom=3cm,left=3cm,right=3cm]{geometry}


\colorlet{titlecolor}{color2}% save the secondary color before redefining it
\definecolor{color2}{rgb}{0.1,0.1,0.1}% redefine the secondary color to purple
\xpatchcmd{\cventry}{.\strut}{\strut}{}{}
% Marge aux 4 coins de la page, ici elles sont réduites pour gagner de la place
%\usepackage[top=1.0cm, bottom=1.0cm, left=1.6cm, right=1.6cm]{geometry}
% Largeur de la colonne de gauche pour les dates
\setlength{\hintscolumnwidth}{2,5cm}
\firstname{Dounia}
\familyname{Darkaoui}
\address{Université de Caen Normandie \\ Laboratoire de Mathématiques Nicolas Oresme \\ UFR des sciences - Campus 2   \\ 14032
Caen, France}{ \\Office: S3 111 \\}
\email{dounia.darkaoui@unicaen.fr}
%\homepage{}
%\mobile{+336 52 39 83 27}
%\extrainfo{Born on 22/02/2000}
\begin{document}
\maketitle
% Marge négative entre le titre et la partie expérience, pour gagner de la place
\section{\textbf{Background}}
\cventry{2024-ongoing}{PhD in Number Theory}{Universit\'e de Caen Normandie (LMNO)}{}{}{Advisors: \textit{Denis Simon} and \textit{Martin Weimann}\\Title: \textit{Applications and generalisations of the OM algorithm for the
factorisation of polynomials in local fields}}

\cventry{2020-2024}{ENS Rennes}{}{}{}{}
\cventry{2021-2024}{Master's degree Mathematics and Applications}{Algebra and Geometry}{Université de Rennes}{}{with highest honor}
\cventry{2023}{Agrégation de Mathématiques}{Option C (Algebra and formal calculus)}{(french teaching exam)}{}{}
\cventry{2021}{Bachelor's degree Mathematics and Applications}{}{Université de Rennes}{}{with highest honor}
\cventry{2018}{Baccalauréat - General Scientific}{Option: Mathematics}{Lycée Guillaume Apollinaire}{Nice, France}{}


\section{\textbf{Experiences}}
\cventry{Apr.-Jul. 2024}{Research Internship}{Universit\'e de Caen Normandie}{France}{}{Advisors: \textit{Prof. Denis Simon} and \textit{MCF. Martin Weimann}\\Title: \textit{On polynomials factoristation in complete fields}}

\cventry{May-June 2022}{Research Internship}{University of Copenhagen}{Denmark}{}{Advisor: \textit{Nathalie Wahl}\\Title: \textit{Lens spaces (algebraic topology)}}

\cventry{Apr.-May 2021}{Research Internship}{Laboratoire J-A Dieudonné}{Nice, France}{}{Advisor: \textit{Indira Chatterji}\\Title: \textit{Polya's theorem and random walks in groups}}

\section{\textbf{Organization}}
\cventry{Dec. 2020}{Organizer}{Rendez-vous des jeunes mathématiciennes et informaticiennes \textbf{RJMI} (Rennes)}{the goal of the event is to support and help female high school students to pursue in science majors such as mathematics and computer science.}{}{}

\section{\textbf{Outreach}}

\cventry{Jan. 2024}{RJMI}{}{Conference on tropical geometry intended for an audience of high school students}{}{}

\cventry{Dec. 2022}{RJMI}{}{Supervised research activity with a group of high school students}{}{}

\cventry{Dec. 2021}{RJMI}{}{Supervised research activity with a group of high school students}{}{}

\section{\textbf{Talks}}
\cvitem{Apr. 2026}{\textbf{Fast computation of Riemann--Roch spaces} (Cheltenham, France), WINGs 2026}
\cvitem{Mars. 2026}{\textbf{Fast computation of Riemann--Roch spaces} (Marseille, France), Journées nationales de calcul formel, CIRM}
\cvitem{Dec. 2025}{\textbf{Espaces de Riemann--Roch : de la géométrie à l'arithmétique} (Lille, France), Young researchers seminar, Laboratoire Paul Painlevé}
\cvitem{Oct. 2025}{\textbf{De la factorisation de polynômes aux espaces de Riemann-Roch} (Marseille, France), Summer school Ethén, CIRM}
\cvitem{May 2025}{\textbf{Espaces de Riemann--Roch et calcul de Base} (Bordeaux, France), Lambda seminar, IMB}
\cvitem{May 2025}{\textbf{Espaces de Riemann--Roch et calcul de Base} (Caen, France), Young researchers seminar, LMNO}




\section{\textbf{Formation}}


\cvitem{May 2025}{\textbf{Formation on discrimination and violence}, ED MIIS (Caen, France)}
\cvitem{Jan. 2025}{\textbf{Scientific Integrity}, MOOC  (Bordeaux, France)}

\cvitem{Dec. 2024}{\textbf{Formation on educational resources}, INSPE - ED MIIS (Caen, France)}



\section{\textbf{Workshops and Conferences}}
\cvitem{Apr. 2026}{\textbf{Women in Number Theory and Geometry  } (Cheltenham, England), Retreat}
\cvitem{Mar. 2026}{\textbf{Journées nationales de calcul formel } (Marseille, France), French Computer Algebra Days}
\cvitem{Jan. 2026}{\textbf{Journées Louis Antoine: Algebraic curves over finite fields and their rational points} (Rennes, France)}
\cvitem{Jan. 2026}{\textbf{Final ANR meeting: Barracuda} (Marseille, France)}
\cvitem{Dec. 2025}{\textbf{Number theory retreat} (Bayeux, France)}
\cvitem{Oct. 2025}{\textbf{EthéN} (Marseille, France), Summer school on number theory}
\cvitem{Jul. 2025}{\textbf{Journées arithmétiques 2025} (Esch-sur-Alzette, Luxembourg)}
\cvitem{Mar. 2025}{\textbf{Journées nationales de calcul formel } (Marseille, France), French Computer Algebra Days}
\cvitem{Feb. 2025}{\textbf{Sagedays 128} (Le Teich, France)}
\cvitem{Jan. 2025}{\textbf{ANR meeting 2025: Barracuda} (Bayeux, France)}
\cvitem{Nov. 2024}{\textbf{Number theory retreat} (Bayeux, France)}
\cvitem{Sep. 2024}{\textbf{EthéN} (Marseille, France), Summer school on number theory}

\section{\textbf{Teaching}}
\cventry{2025}{\textbf{Teaching Assistant}}{}{}{}{ \begin{itemize}[leftmargin=15mm]
\item Algorithmics (second-year-mathematics students) (55h) 
\item Analysis and visualisation on Geogebra (second-year-mathematics students) (9h)
\end{itemize}}
\cventry{2024}{\textbf{Teaching Assistant}}{}{}{}{ \begin{itemize}[leftmargin=15mm]
\item Algorithmics (second-year-mathematics students) (35h) 
\item Python (third-year-mathematics students) (25h)
\end{itemize}}

\cventry{2023-2024}{\textbf{Colles CPGE}}{Lycée Chateaubriand}{Rennes, France}{}
{MP* and MP (second-year-mathematics students) ($\sim$ 45h) }



\section{\textbf{Skills}}
\cvitem{Languages}{\textbf{French} (\textit{native}), \textbf{English} (\textit{fluent}), \textbf{Arabic} (\textit{basics})}
\cvitem{Programming}{ \textbf{SageMath}, \textbf{LaTeX}, \textbf{Python}, \textbf{OCaml} }


%\newpage

%\textbf{\huge{Cover Letter} }
%\\

%Dear organizers,
%\\

%I am a first-year PhD in the Université de Caen Normandie, France. I am working in the number theory team of the Laboratoire de Nicolas Oresme. My advisors are Martin Weimann and Denis Simon. My PhD is on applications and generalizations of the OM algorithm which factors polynomials in local fields. I obtained two Master's degrees, one with a specialization in teaching and another in pure mathematics.  
%My current research topic is on the computation of bases of Riemann-Roch spaces. Therefore, I generally work with valuations, fractional ideals and divisors in function fields. I am looking forward to talk about my research in this summer school as my subject mixes different areas of number theory.
%\\

%Thank you for taking into consideration my application.
%\\

%Best regards,

%Dounia Darkaoui










\end{document}
